% The simulated readers (runs 11 and 12, both side runs). Every number is in a table generated by paper/summary_tables/readers.py, which reads paper/data/run11-simulated-readers and paper/data/run12-simulated-readers-api through explicit paths, whatever run folder the build is given. The design, the three readings and their words were fixed in docs/superpowers/specs/2026-11-18-llm-simulated-readers-design.md and docs/superpowers/plans/2026-11-18-simulated-readers.md before the first call; the second run and the rule for reading it in docs/superpowers/plans/2026-11-19-simulated-readers-api-run.md before any reply was scored.
\section{Simulated Readers and Their Alignment with Earlier Tests}
\label{sec:readers}
% src: docs/superpowers/specs/2026-11-18-llm-simulated-readers-design.md (What had been seen when this was written; Wording the result may take); paper/data/run11-simulated-readers/README.md (a side run; the code state; one outcome figure known beforehand); no number
A language model can be asked to play a reader and say whether that reader would click a headline. Such simulated audiences are offered as a substitute for a live test, and published work finds that their raw answers need calibration against human outcomes \cite{persson2026statistical,ankargren2026when,ashokkumar2026large}. This section asks two questions of the headline archive. The first is whether a panel of simulated readers tells which headline of a test drew more clicks. The second is whether aligning the panel with earlier results helps, when the alignment may learn only from tests that the rule of Section~\ref{sec:setting} makes visible. That rule cannot bind the model itself, whose training ends years after the archive was published. We added this analysis after every other result of the paper was known. Its design, three readings and the words each reading may take were written down before the first call to a model on an archive headline (Appendix~\ref{sec:analysis-history}). It is exploratory, like the rest of the paper outside Section~\ref{sec:confirmatory}, and it has no fresh split to confirm it. Appendix~\ref{sec:readers-detail} holds the prompts, the reader profiles and the remaining measures.

% src: paper/data/run11-simulated-readers/plan.json counts.tests 1770, counts.contrasts 5966, counts.blocks 63, counts.feeds 308, block_size 30, readers 10; tables/fmt-readers-more.tex (units ranked 1,771)
\paragraph{One reader, one feed.} The rows are the main run's scored rows in the headline-only cell, where the arms of a test differ in the headline alone: 1,770 tests with 5,966 contrasts. In the real test a visitor saw one package and clicked or did not. A simulated reader therefore never sees two headlines of one test. One call to the model plays one reader who scrolls one feed of up to 30 story cards. Each card shows the headline of one arm, each from a different test, and nothing else. For each card the reader returns a probability from 0 to 100 that they would click it. The tests of a scored origin are cut into blocks of 30. A block has as many feeds as its largest test has arms, and every feed of a block holds the same stories in the same positions, so that two arms of a test differ only in the headline shown. A test with fewer arms repeats them to fill the feeds, and an arm's score averages its repeats. This gives 63 blocks and 308 feeds.

% src: paper/data/run11-simulated-readers/README.md (the three models, their training data cutoff of June 2026, the primary model declared before the run); src/mootcourt/simulators/readers/protocol.py SYSTEM_PROMPT (between 2013 and 2015; no publisher named), PROFILES (ten), HALVES; no measured number; 5.5 is part of each model's name
\paragraph{The panel.} Ten reader profiles read every feed. A profile is a short paragraph that gives an age band, an occupation, interests, a mood, a device and how the reader met the story. The archive holds no audience data, so the profiles are assumed, and their spread is variation between prompts and not a sample of the real audience. The instructions place the reader between 2013 and 2015 and do not name the publisher. Each call's probabilities are centred on that call's own mean. An arm's \emph{panel score} is the mean of the centred probabilities over the readers. For a contrast, the panel's direction is the sign of the variant's score minus the reference's, and a difference of exactly zero counts as half right. We ran three models of one vendor, from the largest to the smallest Claude Opus 5.5, Claude Sonnet 5.5 and Claude Haiku 5.5, and declared the largest as primary before the run. The archive was public years before the training data of all three ends, so nothing is hidden from them, and the models may have met these headlines. Three probes describe that risk and cannot exclude it (Appendix~\ref{sec:readers-detail}).

% src: tables/fmt-readers-readings.tex rows the earliest-arm rule 0.5701, the text classifier 0.5484, the direction ceiling 0.8301; paper/data/run11-simulated-readers/simulated-readers-claude-opus-5-5-20261106-2dc461f52ea4.json baselines.always_reference.accuracy 0.5701, baselines.text_lr.accuracy 0.5484, baselines.ceiling.accuracy 0.8301
\paragraph{References.} The measure is \emph{direction accuracy}: the share of contrasts on which a method names the sign of the observed effect. Three references are computed on the same rows. The \emph{earliest-arm rule} always says that the variant loses to the reference arm, the package created first, and it is right on 0.570 of the contrasts. The text classifier of Section~\ref{sec:method}, fitted on earlier tests, is right on 0.548. The ceiling of Section~\ref{sec:protocol}, an estimate of the accuracy a method would reach if it knew each test's true effect, is 0.830. A rule that always picks the first headline written thus scores 0.570 without reading anything, so one half is not the reference to beat.

% src: docs/superpowers/specs/2026-11-18-llm-simulated-readers-design.md, Aligning the readers with earlier results (seven methods; refitted at every scored origin on the tests the visibility rule lets a forecast see; ten earlier tests in the prompt; 200 earlier tests and a note of at most 150 words); src/mootcourt/simulators/readers/protocol.py N_EXAMPLES 10, N_BRIEF_UNITS 200, NOTE_WORDS 150
\paragraph{Seven ways to align the panel.} Each method is refitted at every scored origin and may use only the headline-only tests that are visible at that origin under the rule of Section~\ref{sec:setting}. At the first scored origin no such test exists, so every aligned panel equals the plain panel there. No prompt, profile or setting was tuned on an outcome. Five methods use the stored answers of the plain panel.
\begin{enumerate}
\item \emph{Direction map.} A logistic regression of the observed direction on the panel's difference.
\item \emph{Reader weights.} Weights over the ten readers, fitted so that the weighted panel explains the click counts of earlier arms.
\item \emph{Audience matching.} Weights chosen only so that the panel's mean level matches the mean click rate of earlier arms, with no use of which arm won.
\item \emph{Error correction.} A monotone map from panel score to click rate, then a linear correction on four fixed traits of a headline.
\item \emph{Text classifier with the panel.} The panel's difference enters the text classifier as one more feature.
\end{enumerate}
Two methods put earlier results in front of the reader and need new calls, on the feeds of the plain panel.
\begin{enumerate}\setcounter{enumi}{5}
\item \emph{Earlier results in the prompt.} Before the feed, the reader sees ten earlier tests with the click rate of each headline.
\item \emph{Audience brief.} The model reads 200 earlier tests with their click rates and writes a note of at most 150 words on what this audience clicked. Every reader of that origin receives the note.
\end{enumerate}

% src: docs/superpowers/specs/2026-11-18-llm-simulated-readers-design.md, Measures (three readings; a 90 percent bootstrap interval from 500 resamples of whole blocks; Bonferroni-adjusted bounds for a family of three; the adjusted interval governs the wording); paper/data/run11-simulated-readers/simulated-readers-claude-opus-5-5-20261106-2dc461f52ea4.json readings.weights_minus_plain.family_size 3, readings.weights_minus_plain.resamples 500, counts.blocks 63
\paragraph{Three readings, fixed beforehand.} Each reading is a paired difference in direction accuracy on all contrasts: (1) the plain panel minus the earliest-arm rule, (2) the panel with reader weights minus the plain panel, and (3) the text classifier with the panel minus the text classifier alone. Each carries a 90 percent bootstrap interval from 500 resamples of whole blocks, with bounds adjusted for the family of three readings as in Section~\ref{sec:protocol}. For the primary model the adjusted interval chooses among three sentences fixed beforehand, for an interval above zero, containing zero and below zero. The two smaller models and every other comparison are described without a verdict word.

% src: paper/data/run11-simulated-readers/README.md (the transport; what the tool adds: the present date beside the instruction to play a person of 2013 to 2015; no control of sampling); paper/data/run12-simulated-readers-api/README.md (the same 9,007 requests per model through the vendor's API; four request fields; no sampling setting can be sent; planned before any reply of run 11 was scored); docs/superpowers/plans/2026-11-19-simulated-readers-api-run.md, The comparison between the two runs, fixed now; tables/fmt-readers-more.tex row planned 9,007; run 11 and run 12 are the names of the two run folders
\paragraph{Two runs of the same requests.} The first run (run 11) reached the models through the vendor's command line tool. That tool puts its own context around each call, the present date included, so a reader was told the year 2026 beside the instruction to play a person of 2013 to 2015. We therefore planned a second run (run 12) before any reply of the first was scored. It repeats the same 9,007 planned calls per model through the vendor's programming interface, where the model receives our two texts and nothing else. The reader calls of method 7 carry the note written in that run, so their text differs between the runs, and every other request is identical. Neither route lets us fix the sampling of a reply, so two runs of one request differ by chance as well. The rule for reading the second run was fixed with its plan: if, for the primary model, the interval of the difference in the plain panel's accuracy between the runs contains zero, we say that the added context did not move that accuracy by an amount this design can detect. The fixed words of the three readings come from run 11 in either case.

\begin{table*}[t]\centering\footnotesize
% src: tables/fmt-readers-readings.tex; paper/data/run11-simulated-readers/simulated-readers-claude-opus-5-5-20261106-2dc461f52ea4.json counts.contrasts 5966, counts.blocks 63, readings.plain_minus_always_reference.family_size 3, readings.plain_minus_always_reference.resamples 500, readings.plain_minus_always_reference.verdict contains_zero, readings.weights_minus_plain.verdict above, readings.text_with_panel_minus_text.verdict above
\caption{The three readings of the simulated readers for the three models, in the run through the command line tool (run 11) and in the run of the same requests through the vendor's programming interface (run 12). The upper rows give the direction accuracy of the three references and of each model's plain panel on the 5,966 contrasts of the headline-only cell. Each reading is a paired difference in direction accuracy with an interval at 90 percent from a bootstrap of 500 resamples of whole blocks. The bounds printed are the ones adjusted for the family of three readings. The last column gives the words fixed before the first call, chosen by the adjusted interval of the primary model in run 11. No other row carries a verdict.}
\label{tab:fmt-readers-readings}
% generated by paper/build_tables.py from the report keys named above each row (readers); it reads paper/data/run11-simulated-readers and paper/data/run12-simulated-readers-api
\begin{tabular}{p{4.9cm}rrp{2.5cm}}
\toprule
reading and model & run 11 & run 12 & fixed words, \\
 & & & Opus in run 11 \\
\midrule
\multicolumn{4}{@{}l}{\emph{Share of contrasts whose direction is matched}} \\
% src: paper/data/run11-simulated-readers/simulated-readers-claude-opus-5-5-20261106-2dc461f52ea4.json baselines.always_reference.accuracy; paper/data/run12-simulated-readers-api/simulated-readers-claude-opus-5-5-run12-api-20261106-2dc461f52ea4.json baselines.always_reference.accuracy
\quad the earliest-arm rule & 0.5701 & 0.5701 &  \\
% src: paper/data/run11-simulated-readers/simulated-readers-claude-opus-5-5-20261106-2dc461f52ea4.json baselines.text_lr.accuracy; paper/data/run12-simulated-readers-api/simulated-readers-claude-opus-5-5-run12-api-20261106-2dc461f52ea4.json baselines.text_lr.accuracy
\quad the text classifier & 0.5484 & 0.5484 &  \\
% src: paper/data/run11-simulated-readers/simulated-readers-claude-opus-5-5-20261106-2dc461f52ea4.json baselines.ceiling.accuracy; paper/data/run12-simulated-readers-api/simulated-readers-claude-opus-5-5-run12-api-20261106-2dc461f52ea4.json baselines.ceiling.accuracy
\quad the direction ceiling & 0.8301 & 0.8301 &  \\
% src: paper/data/run11-simulated-readers/simulated-readers-claude-opus-5-5-20261106-2dc461f52ea4.json plain.accuracy; paper/data/run12-simulated-readers-api/simulated-readers-claude-opus-5-5-run12-api-20261106-2dc461f52ea4.json plain.accuracy
\quad plain panel, Opus & 0.5742 & 0.5742 &  \\
% src: paper/data/run11-simulated-readers/simulated-readers-claude-sonnet-5-5-20261106-2dc461f52ea4.json plain.accuracy; paper/data/run12-simulated-readers-api/simulated-readers-claude-sonnet-5-5-run12-api-20261106-2dc461f52ea4.json plain.accuracy
\quad plain panel, Sonnet & 0.5234 & 0.5279 &  \\
% src: paper/data/run11-simulated-readers/simulated-readers-claude-haiku-5-5-20261106-2dc461f52ea4.json plain.accuracy; paper/data/run12-simulated-readers-api/simulated-readers-claude-haiku-5-5-run12-api-20261106-2dc461f52ea4.json plain.accuracy
\quad plain panel, Haiku & 0.5144 & 0.5223 &  \\
\midrule
\multicolumn{4}{@{}l}{\emph{1. Plain panel minus the earliest-arm rule}} \\
% src: paper/data/run11-simulated-readers/simulated-readers-claude-opus-5-5-20261106-2dc461f52ea4.json readings.plain_minus_always_reference.estimate; paper/data/run11-simulated-readers/simulated-readers-claude-opus-5-5-20261106-2dc461f52ea4.json readings.plain_minus_always_reference.adjusted_lo; paper/data/run11-simulated-readers/simulated-readers-claude-opus-5-5-20261106-2dc461f52ea4.json readings.plain_minus_always_reference.adjusted_hi; paper/data/run12-simulated-readers-api/simulated-readers-claude-opus-5-5-run12-api-20261106-2dc461f52ea4.json readings.plain_minus_always_reference.estimate; paper/data/run12-simulated-readers-api/simulated-readers-claude-opus-5-5-run12-api-20261106-2dc461f52ea4.json readings.plain_minus_always_reference.adjusted_lo; paper/data/run12-simulated-readers-api/simulated-readers-claude-opus-5-5-run12-api-20261106-2dc461f52ea4.json readings.plain_minus_always_reference.adjusted_hi; paper/data/run11-simulated-readers/simulated-readers-claude-opus-5-5-20261106-2dc461f52ea4.json readings.plain_minus_always_reference.adjusted_lo; paper/data/run11-simulated-readers/simulated-readers-claude-opus-5-5-20261106-2dc461f52ea4.json readings.plain_minus_always_reference.adjusted_hi; paper/data/run11-simulated-readers/simulated-readers-claude-opus-5-5-20261106-2dc461f52ea4.json readings.plain_minus_always_reference.verdict
\quad Opus & 0.0041 (-0.0174, 0.0256) & 0.0041 (-0.0171, 0.0253) & not shown to match more often \\
% src: paper/data/run11-simulated-readers/simulated-readers-claude-sonnet-5-5-20261106-2dc461f52ea4.json readings.plain_minus_always_reference.estimate; paper/data/run11-simulated-readers/simulated-readers-claude-sonnet-5-5-20261106-2dc461f52ea4.json readings.plain_minus_always_reference.adjusted_lo; paper/data/run11-simulated-readers/simulated-readers-claude-sonnet-5-5-20261106-2dc461f52ea4.json readings.plain_minus_always_reference.adjusted_hi; paper/data/run12-simulated-readers-api/simulated-readers-claude-sonnet-5-5-run12-api-20261106-2dc461f52ea4.json readings.plain_minus_always_reference.estimate; paper/data/run12-simulated-readers-api/simulated-readers-claude-sonnet-5-5-run12-api-20261106-2dc461f52ea4.json readings.plain_minus_always_reference.adjusted_lo; paper/data/run12-simulated-readers-api/simulated-readers-claude-sonnet-5-5-run12-api-20261106-2dc461f52ea4.json readings.plain_minus_always_reference.adjusted_hi
\quad Sonnet & -0.0467 (-0.0704, -0.0230) & -0.0422 (-0.0637, -0.0206) &  \\
% src: paper/data/run11-simulated-readers/simulated-readers-claude-haiku-5-5-20261106-2dc461f52ea4.json readings.plain_minus_always_reference.estimate; paper/data/run11-simulated-readers/simulated-readers-claude-haiku-5-5-20261106-2dc461f52ea4.json readings.plain_minus_always_reference.adjusted_lo; paper/data/run11-simulated-readers/simulated-readers-claude-haiku-5-5-20261106-2dc461f52ea4.json readings.plain_minus_always_reference.adjusted_hi; paper/data/run12-simulated-readers-api/simulated-readers-claude-haiku-5-5-run12-api-20261106-2dc461f52ea4.json readings.plain_minus_always_reference.estimate; paper/data/run12-simulated-readers-api/simulated-readers-claude-haiku-5-5-run12-api-20261106-2dc461f52ea4.json readings.plain_minus_always_reference.adjusted_lo; paper/data/run12-simulated-readers-api/simulated-readers-claude-haiku-5-5-run12-api-20261106-2dc461f52ea4.json readings.plain_minus_always_reference.adjusted_hi
\quad Haiku & -0.0556 (-0.0791, -0.0322) & -0.0478 (-0.0707, -0.0248) &  \\
\midrule
\multicolumn{4}{@{}l}{\emph{2. Panel with reader weights minus the plain panel}} \\
% src: paper/data/run11-simulated-readers/simulated-readers-claude-opus-5-5-20261106-2dc461f52ea4.json readings.weights_minus_plain.estimate; paper/data/run11-simulated-readers/simulated-readers-claude-opus-5-5-20261106-2dc461f52ea4.json readings.weights_minus_plain.adjusted_lo; paper/data/run11-simulated-readers/simulated-readers-claude-opus-5-5-20261106-2dc461f52ea4.json readings.weights_minus_plain.adjusted_hi; paper/data/run12-simulated-readers-api/simulated-readers-claude-opus-5-5-run12-api-20261106-2dc461f52ea4.json readings.weights_minus_plain.estimate; paper/data/run12-simulated-readers-api/simulated-readers-claude-opus-5-5-run12-api-20261106-2dc461f52ea4.json readings.weights_minus_plain.adjusted_lo; paper/data/run12-simulated-readers-api/simulated-readers-claude-opus-5-5-run12-api-20261106-2dc461f52ea4.json readings.weights_minus_plain.adjusted_hi; paper/data/run11-simulated-readers/simulated-readers-claude-opus-5-5-20261106-2dc461f52ea4.json readings.weights_minus_plain.adjusted_lo; paper/data/run11-simulated-readers/simulated-readers-claude-opus-5-5-20261106-2dc461f52ea4.json readings.weights_minus_plain.adjusted_hi; paper/data/run11-simulated-readers/simulated-readers-claude-opus-5-5-20261106-2dc461f52ea4.json readings.weights_minus_plain.verdict
\quad Opus & 0.0140 (0.0004, 0.0276) & 0.0110 (-0.0035, 0.0254) & raised \\
% src: paper/data/run11-simulated-readers/simulated-readers-claude-sonnet-5-5-20261106-2dc461f52ea4.json readings.weights_minus_plain.estimate; paper/data/run11-simulated-readers/simulated-readers-claude-sonnet-5-5-20261106-2dc461f52ea4.json readings.weights_minus_plain.adjusted_lo; paper/data/run11-simulated-readers/simulated-readers-claude-sonnet-5-5-20261106-2dc461f52ea4.json readings.weights_minus_plain.adjusted_hi; paper/data/run12-simulated-readers-api/simulated-readers-claude-sonnet-5-5-run12-api-20261106-2dc461f52ea4.json readings.weights_minus_plain.estimate; paper/data/run12-simulated-readers-api/simulated-readers-claude-sonnet-5-5-run12-api-20261106-2dc461f52ea4.json readings.weights_minus_plain.adjusted_lo; paper/data/run12-simulated-readers-api/simulated-readers-claude-sonnet-5-5-run12-api-20261106-2dc461f52ea4.json readings.weights_minus_plain.adjusted_hi
\quad Sonnet & 0.0365 (0.0222, 0.0507) & 0.0326 (0.0162, 0.0490) &  \\
% src: paper/data/run11-simulated-readers/simulated-readers-claude-haiku-5-5-20261106-2dc461f52ea4.json readings.weights_minus_plain.estimate; paper/data/run11-simulated-readers/simulated-readers-claude-haiku-5-5-20261106-2dc461f52ea4.json readings.weights_minus_plain.adjusted_lo; paper/data/run11-simulated-readers/simulated-readers-claude-haiku-5-5-20261106-2dc461f52ea4.json readings.weights_minus_plain.adjusted_hi; paper/data/run12-simulated-readers-api/simulated-readers-claude-haiku-5-5-run12-api-20261106-2dc461f52ea4.json readings.weights_minus_plain.estimate; paper/data/run12-simulated-readers-api/simulated-readers-claude-haiku-5-5-run12-api-20261106-2dc461f52ea4.json readings.weights_minus_plain.adjusted_lo; paper/data/run12-simulated-readers-api/simulated-readers-claude-haiku-5-5-run12-api-20261106-2dc461f52ea4.json readings.weights_minus_plain.adjusted_hi
\quad Haiku & 0.0231 (0.0073, 0.0389) & -0.0012 (-0.0176, 0.0152) &  \\
\midrule
\multicolumn{4}{@{}l}{\emph{3. Text classifier with the panel minus the text classifier alone}} \\
% src: paper/data/run11-simulated-readers/simulated-readers-claude-opus-5-5-20261106-2dc461f52ea4.json readings.text_with_panel_minus_text.estimate; paper/data/run11-simulated-readers/simulated-readers-claude-opus-5-5-20261106-2dc461f52ea4.json readings.text_with_panel_minus_text.adjusted_lo; paper/data/run11-simulated-readers/simulated-readers-claude-opus-5-5-20261106-2dc461f52ea4.json readings.text_with_panel_minus_text.adjusted_hi; paper/data/run12-simulated-readers-api/simulated-readers-claude-opus-5-5-run12-api-20261106-2dc461f52ea4.json readings.text_with_panel_minus_text.estimate; paper/data/run12-simulated-readers-api/simulated-readers-claude-opus-5-5-run12-api-20261106-2dc461f52ea4.json readings.text_with_panel_minus_text.adjusted_lo; paper/data/run12-simulated-readers-api/simulated-readers-claude-opus-5-5-run12-api-20261106-2dc461f52ea4.json readings.text_with_panel_minus_text.adjusted_hi; paper/data/run11-simulated-readers/simulated-readers-claude-opus-5-5-20261106-2dc461f52ea4.json readings.text_with_panel_minus_text.adjusted_lo; paper/data/run11-simulated-readers/simulated-readers-claude-opus-5-5-20261106-2dc461f52ea4.json readings.text_with_panel_minus_text.adjusted_hi; paper/data/run11-simulated-readers/simulated-readers-claude-opus-5-5-20261106-2dc461f52ea4.json readings.text_with_panel_minus_text.verdict
\quad Opus & 0.0203 (0.0100, 0.0306) & 0.0184 (0.0092, 0.0276) & raised \\
% src: paper/data/run11-simulated-readers/simulated-readers-claude-sonnet-5-5-20261106-2dc461f52ea4.json readings.text_with_panel_minus_text.estimate; paper/data/run11-simulated-readers/simulated-readers-claude-sonnet-5-5-20261106-2dc461f52ea4.json readings.text_with_panel_minus_text.adjusted_lo; paper/data/run11-simulated-readers/simulated-readers-claude-sonnet-5-5-20261106-2dc461f52ea4.json readings.text_with_panel_minus_text.adjusted_hi; paper/data/run12-simulated-readers-api/simulated-readers-claude-sonnet-5-5-run12-api-20261106-2dc461f52ea4.json readings.text_with_panel_minus_text.estimate; paper/data/run12-simulated-readers-api/simulated-readers-claude-sonnet-5-5-run12-api-20261106-2dc461f52ea4.json readings.text_with_panel_minus_text.adjusted_lo; paper/data/run12-simulated-readers-api/simulated-readers-claude-sonnet-5-5-run12-api-20261106-2dc461f52ea4.json readings.text_with_panel_minus_text.adjusted_hi
\quad Sonnet & 0.0000 (-0.0068, 0.0068) & -0.0007 (-0.0069, 0.0056) &  \\
% src: paper/data/run11-simulated-readers/simulated-readers-claude-haiku-5-5-20261106-2dc461f52ea4.json readings.text_with_panel_minus_text.estimate; paper/data/run11-simulated-readers/simulated-readers-claude-haiku-5-5-20261106-2dc461f52ea4.json readings.text_with_panel_minus_text.adjusted_lo; paper/data/run11-simulated-readers/simulated-readers-claude-haiku-5-5-20261106-2dc461f52ea4.json readings.text_with_panel_minus_text.adjusted_hi; paper/data/run12-simulated-readers-api/simulated-readers-claude-haiku-5-5-run12-api-20261106-2dc461f52ea4.json readings.text_with_panel_minus_text.estimate; paper/data/run12-simulated-readers-api/simulated-readers-claude-haiku-5-5-run12-api-20261106-2dc461f52ea4.json readings.text_with_panel_minus_text.adjusted_lo; paper/data/run12-simulated-readers-api/simulated-readers-claude-haiku-5-5-run12-api-20261106-2dc461f52ea4.json readings.text_with_panel_minus_text.adjusted_hi
\quad Haiku & 0.0015 (-0.0045, 0.0075) & 0.0052 (-0.0017, 0.0121) &  \\
\bottomrule
\end{tabular}

\end{table*}

\begin{table*}[t]\centering\footnotesize
% src: tables/fmt-readers-methods.tex; paper/data/run11-simulated-readers/simulated-readers-claude-opus-5-5-20261106-2dc461f52ea4.json methods.audience_brief.minus_always_reference.resamples 500, counts.blocks 63; run 11 and run 12 are the names of the two run folders
\caption{The plain panel and the seven alignment methods in run 11, by model. The first column is direction accuracy. The other two are paired differences from the plain panel and from the earliest-arm rule, each with an interval at 90 percent from a bootstrap of 500 resamples of whole blocks. These intervals belong to no declared family and are not adjusted for the number of comparisons in the table, so they describe and carry no verdict. Table~\ref{tab:fmt-readers-methods-api} is the same table for run 12.}
\label{tab:fmt-readers-methods}
% generated by paper/build_tables.py from the report keys named above each row (readers); it reads paper/data/run11-simulated-readers
\begin{tabular}{p{4.6cm}rrr}
\toprule
model and method & share matched & minus the plain panel & minus the earliest-arm rule \\
\midrule
\multicolumn{4}{@{}l}{\emph{Opus}} \\
% src: paper/data/run11-simulated-readers/simulated-readers-claude-opus-5-5-20261106-2dc461f52ea4.json plain.accuracy; paper/data/run11-simulated-readers/simulated-readers-claude-opus-5-5-20261106-2dc461f52ea4.json readings.plain_minus_always_reference.estimate; paper/data/run11-simulated-readers/simulated-readers-claude-opus-5-5-20261106-2dc461f52ea4.json readings.plain_minus_always_reference.lo; paper/data/run11-simulated-readers/simulated-readers-claude-opus-5-5-20261106-2dc461f52ea4.json readings.plain_minus_always_reference.hi
\quad plain panel & 0.5742 &  & 0.0041 (-0.0117, 0.0217) \\
% src: paper/data/run11-simulated-readers/simulated-readers-claude-opus-5-5-20261106-2dc461f52ea4.json methods.direction_map.accuracy; paper/data/run11-simulated-readers/simulated-readers-claude-opus-5-5-20261106-2dc461f52ea4.json methods.direction_map.minus_plain.estimate; paper/data/run11-simulated-readers/simulated-readers-claude-opus-5-5-20261106-2dc461f52ea4.json methods.direction_map.minus_plain.lo; paper/data/run11-simulated-readers/simulated-readers-claude-opus-5-5-20261106-2dc461f52ea4.json methods.direction_map.minus_plain.hi; paper/data/run11-simulated-readers/simulated-readers-claude-opus-5-5-20261106-2dc461f52ea4.json methods.direction_map.minus_always_reference.estimate; paper/data/run11-simulated-readers/simulated-readers-claude-opus-5-5-20261106-2dc461f52ea4.json methods.direction_map.minus_always_reference.lo; paper/data/run11-simulated-readers/simulated-readers-claude-opus-5-5-20261106-2dc461f52ea4.json methods.direction_map.minus_always_reference.hi
\quad 1. direction map & 0.5840 & 0.0098 (-0.0013, 0.0200) & 0.0139 (0.0039, 0.0247) \\
% src: paper/data/run11-simulated-readers/simulated-readers-claude-opus-5-5-20261106-2dc461f52ea4.json methods.reader_weights.accuracy; paper/data/run11-simulated-readers/simulated-readers-claude-opus-5-5-20261106-2dc461f52ea4.json methods.reader_weights.minus_plain.estimate; paper/data/run11-simulated-readers/simulated-readers-claude-opus-5-5-20261106-2dc461f52ea4.json methods.reader_weights.minus_plain.lo; paper/data/run11-simulated-readers/simulated-readers-claude-opus-5-5-20261106-2dc461f52ea4.json methods.reader_weights.minus_plain.hi; paper/data/run11-simulated-readers/simulated-readers-claude-opus-5-5-20261106-2dc461f52ea4.json methods.reader_weights.minus_always_reference.estimate; paper/data/run11-simulated-readers/simulated-readers-claude-opus-5-5-20261106-2dc461f52ea4.json methods.reader_weights.minus_always_reference.lo; paper/data/run11-simulated-readers/simulated-readers-claude-opus-5-5-20261106-2dc461f52ea4.json methods.reader_weights.minus_always_reference.hi
\quad 2. reader weights & 0.5882 & 0.0140 (0.0041, 0.0241) & 0.0181 (0.0002, 0.0368) \\
% src: paper/data/run11-simulated-readers/simulated-readers-claude-opus-5-5-20261106-2dc461f52ea4.json methods.audience_matching.accuracy; paper/data/run11-simulated-readers/simulated-readers-claude-opus-5-5-20261106-2dc461f52ea4.json methods.audience_matching.minus_plain.estimate; paper/data/run11-simulated-readers/simulated-readers-claude-opus-5-5-20261106-2dc461f52ea4.json methods.audience_matching.minus_plain.lo; paper/data/run11-simulated-readers/simulated-readers-claude-opus-5-5-20261106-2dc461f52ea4.json methods.audience_matching.minus_plain.hi; paper/data/run11-simulated-readers/simulated-readers-claude-opus-5-5-20261106-2dc461f52ea4.json methods.audience_matching.minus_always_reference.estimate; paper/data/run11-simulated-readers/simulated-readers-claude-opus-5-5-20261106-2dc461f52ea4.json methods.audience_matching.minus_always_reference.lo; paper/data/run11-simulated-readers/simulated-readers-claude-opus-5-5-20261106-2dc461f52ea4.json methods.audience_matching.minus_always_reference.hi
\quad 3. audience matching & 0.5716 & -0.0026 (-0.0139, 0.0066) & 0.0015 (-0.0165, 0.0218) \\
% src: paper/data/run11-simulated-readers/simulated-readers-claude-opus-5-5-20261106-2dc461f52ea4.json methods.error_correction.accuracy; paper/data/run11-simulated-readers/simulated-readers-claude-opus-5-5-20261106-2dc461f52ea4.json methods.error_correction.minus_plain.estimate; paper/data/run11-simulated-readers/simulated-readers-claude-opus-5-5-20261106-2dc461f52ea4.json methods.error_correction.minus_plain.lo; paper/data/run11-simulated-readers/simulated-readers-claude-opus-5-5-20261106-2dc461f52ea4.json methods.error_correction.minus_plain.hi; paper/data/run11-simulated-readers/simulated-readers-claude-opus-5-5-20261106-2dc461f52ea4.json methods.error_correction.minus_always_reference.estimate; paper/data/run11-simulated-readers/simulated-readers-claude-opus-5-5-20261106-2dc461f52ea4.json methods.error_correction.minus_always_reference.lo; paper/data/run11-simulated-readers/simulated-readers-claude-opus-5-5-20261106-2dc461f52ea4.json methods.error_correction.minus_always_reference.hi
\quad 4. error correction & 0.5819 & 0.0077 (-0.0016, 0.0157) & 0.0118 (-0.0062, 0.0312) \\
% src: paper/data/run11-simulated-readers/simulated-readers-claude-opus-5-5-20261106-2dc461f52ea4.json methods.text_with_panel.accuracy; paper/data/run11-simulated-readers/simulated-readers-claude-opus-5-5-20261106-2dc461f52ea4.json methods.text_with_panel.minus_plain.estimate; paper/data/run11-simulated-readers/simulated-readers-claude-opus-5-5-20261106-2dc461f52ea4.json methods.text_with_panel.minus_plain.lo; paper/data/run11-simulated-readers/simulated-readers-claude-opus-5-5-20261106-2dc461f52ea4.json methods.text_with_panel.minus_plain.hi; paper/data/run11-simulated-readers/simulated-readers-claude-opus-5-5-20261106-2dc461f52ea4.json methods.text_with_panel.minus_always_reference.estimate; paper/data/run11-simulated-readers/simulated-readers-claude-opus-5-5-20261106-2dc461f52ea4.json methods.text_with_panel.minus_always_reference.lo; paper/data/run11-simulated-readers/simulated-readers-claude-opus-5-5-20261106-2dc461f52ea4.json methods.text_with_panel.minus_always_reference.hi
\quad 5. text classifier with the panel & 0.5687 & -0.0054 (-0.0195, 0.0077) & -0.0013 (-0.0153, 0.0122) \\
% src: paper/data/run11-simulated-readers/simulated-readers-claude-opus-5-5-20261106-2dc461f52ea4.json methods.examples_in_prompt.accuracy; paper/data/run11-simulated-readers/simulated-readers-claude-opus-5-5-20261106-2dc461f52ea4.json methods.examples_in_prompt.minus_plain.estimate; paper/data/run11-simulated-readers/simulated-readers-claude-opus-5-5-20261106-2dc461f52ea4.json methods.examples_in_prompt.minus_plain.lo; paper/data/run11-simulated-readers/simulated-readers-claude-opus-5-5-20261106-2dc461f52ea4.json methods.examples_in_prompt.minus_plain.hi; paper/data/run11-simulated-readers/simulated-readers-claude-opus-5-5-20261106-2dc461f52ea4.json methods.examples_in_prompt.minus_always_reference.estimate; paper/data/run11-simulated-readers/simulated-readers-claude-opus-5-5-20261106-2dc461f52ea4.json methods.examples_in_prompt.minus_always_reference.lo; paper/data/run11-simulated-readers/simulated-readers-claude-opus-5-5-20261106-2dc461f52ea4.json methods.examples_in_prompt.minus_always_reference.hi
\quad 6. earlier results in the prompt & 0.5853 & 0.0111 (0.0019, 0.0217) & 0.0153 (-0.0022, 0.0345) \\
% src: paper/data/run11-simulated-readers/simulated-readers-claude-opus-5-5-20261106-2dc461f52ea4.json methods.audience_brief.accuracy; paper/data/run11-simulated-readers/simulated-readers-claude-opus-5-5-20261106-2dc461f52ea4.json methods.audience_brief.minus_plain.estimate; paper/data/run11-simulated-readers/simulated-readers-claude-opus-5-5-20261106-2dc461f52ea4.json methods.audience_brief.minus_plain.lo; paper/data/run11-simulated-readers/simulated-readers-claude-opus-5-5-20261106-2dc461f52ea4.json methods.audience_brief.minus_plain.hi; paper/data/run11-simulated-readers/simulated-readers-claude-opus-5-5-20261106-2dc461f52ea4.json methods.audience_brief.minus_always_reference.estimate; paper/data/run11-simulated-readers/simulated-readers-claude-opus-5-5-20261106-2dc461f52ea4.json methods.audience_brief.minus_always_reference.lo; paper/data/run11-simulated-readers/simulated-readers-claude-opus-5-5-20261106-2dc461f52ea4.json methods.audience_brief.minus_always_reference.hi
\quad 7. audience brief & 0.6122 & 0.0380 (0.0274, 0.0480) & 0.0422 (0.0260, 0.0575) \\
\midrule
\multicolumn{4}{@{}l}{\emph{Sonnet}} \\
% src: paper/data/run11-simulated-readers/simulated-readers-claude-sonnet-5-5-20261106-2dc461f52ea4.json plain.accuracy; paper/data/run11-simulated-readers/simulated-readers-claude-sonnet-5-5-20261106-2dc461f52ea4.json readings.plain_minus_always_reference.estimate; paper/data/run11-simulated-readers/simulated-readers-claude-sonnet-5-5-20261106-2dc461f52ea4.json readings.plain_minus_always_reference.lo; paper/data/run11-simulated-readers/simulated-readers-claude-sonnet-5-5-20261106-2dc461f52ea4.json readings.plain_minus_always_reference.hi
\quad plain panel & 0.5234 &  & -0.0467 (-0.0648, -0.0275) \\
% src: paper/data/run11-simulated-readers/simulated-readers-claude-sonnet-5-5-20261106-2dc461f52ea4.json methods.direction_map.accuracy; paper/data/run11-simulated-readers/simulated-readers-claude-sonnet-5-5-20261106-2dc461f52ea4.json methods.direction_map.minus_plain.estimate; paper/data/run11-simulated-readers/simulated-readers-claude-sonnet-5-5-20261106-2dc461f52ea4.json methods.direction_map.minus_plain.lo; paper/data/run11-simulated-readers/simulated-readers-claude-sonnet-5-5-20261106-2dc461f52ea4.json methods.direction_map.minus_plain.hi; paper/data/run11-simulated-readers/simulated-readers-claude-sonnet-5-5-20261106-2dc461f52ea4.json methods.direction_map.minus_always_reference.estimate; paper/data/run11-simulated-readers/simulated-readers-claude-sonnet-5-5-20261106-2dc461f52ea4.json methods.direction_map.minus_always_reference.lo; paper/data/run11-simulated-readers/simulated-readers-claude-sonnet-5-5-20261106-2dc461f52ea4.json methods.direction_map.minus_always_reference.hi
\quad 1. direction map & 0.5622 & 0.0388 (0.0233, 0.0550) & -0.0079 (-0.0143, -0.0016) \\
% src: paper/data/run11-simulated-readers/simulated-readers-claude-sonnet-5-5-20261106-2dc461f52ea4.json methods.reader_weights.accuracy; paper/data/run11-simulated-readers/simulated-readers-claude-sonnet-5-5-20261106-2dc461f52ea4.json methods.reader_weights.minus_plain.estimate; paper/data/run11-simulated-readers/simulated-readers-claude-sonnet-5-5-20261106-2dc461f52ea4.json methods.reader_weights.minus_plain.lo; paper/data/run11-simulated-readers/simulated-readers-claude-sonnet-5-5-20261106-2dc461f52ea4.json methods.reader_weights.minus_plain.hi; paper/data/run11-simulated-readers/simulated-readers-claude-sonnet-5-5-20261106-2dc461f52ea4.json methods.reader_weights.minus_always_reference.estimate; paper/data/run11-simulated-readers/simulated-readers-claude-sonnet-5-5-20261106-2dc461f52ea4.json methods.reader_weights.minus_always_reference.lo; paper/data/run11-simulated-readers/simulated-readers-claude-sonnet-5-5-20261106-2dc461f52ea4.json methods.reader_weights.minus_always_reference.hi
\quad 2. reader weights & 0.5598 & 0.0365 (0.0253, 0.0474) & -0.0102 (-0.0308, 0.0082) \\
% src: paper/data/run11-simulated-readers/simulated-readers-claude-sonnet-5-5-20261106-2dc461f52ea4.json methods.audience_matching.accuracy; paper/data/run11-simulated-readers/simulated-readers-claude-sonnet-5-5-20261106-2dc461f52ea4.json methods.audience_matching.minus_plain.estimate; paper/data/run11-simulated-readers/simulated-readers-claude-sonnet-5-5-20261106-2dc461f52ea4.json methods.audience_matching.minus_plain.lo; paper/data/run11-simulated-readers/simulated-readers-claude-sonnet-5-5-20261106-2dc461f52ea4.json methods.audience_matching.minus_plain.hi; paper/data/run11-simulated-readers/simulated-readers-claude-sonnet-5-5-20261106-2dc461f52ea4.json methods.audience_matching.minus_always_reference.estimate; paper/data/run11-simulated-readers/simulated-readers-claude-sonnet-5-5-20261106-2dc461f52ea4.json methods.audience_matching.minus_always_reference.lo; paper/data/run11-simulated-readers/simulated-readers-claude-sonnet-5-5-20261106-2dc461f52ea4.json methods.audience_matching.minus_always_reference.hi
\quad 3. audience matching & 0.5357 & 0.0123 (0.0031, 0.0211) & -0.0344 (-0.0515, -0.0189) \\
% src: paper/data/run11-simulated-readers/simulated-readers-claude-sonnet-5-5-20261106-2dc461f52ea4.json methods.error_correction.accuracy; paper/data/run11-simulated-readers/simulated-readers-claude-sonnet-5-5-20261106-2dc461f52ea4.json methods.error_correction.minus_plain.estimate; paper/data/run11-simulated-readers/simulated-readers-claude-sonnet-5-5-20261106-2dc461f52ea4.json methods.error_correction.minus_plain.lo; paper/data/run11-simulated-readers/simulated-readers-claude-sonnet-5-5-20261106-2dc461f52ea4.json methods.error_correction.minus_plain.hi; paper/data/run11-simulated-readers/simulated-readers-claude-sonnet-5-5-20261106-2dc461f52ea4.json methods.error_correction.minus_always_reference.estimate; paper/data/run11-simulated-readers/simulated-readers-claude-sonnet-5-5-20261106-2dc461f52ea4.json methods.error_correction.minus_always_reference.lo; paper/data/run11-simulated-readers/simulated-readers-claude-sonnet-5-5-20261106-2dc461f52ea4.json methods.error_correction.minus_always_reference.hi
\quad 4. error correction & 0.5484 & 0.0251 (0.0169, 0.0329) & -0.0216 (-0.0404, -0.0042) \\
% src: paper/data/run11-simulated-readers/simulated-readers-claude-sonnet-5-5-20261106-2dc461f52ea4.json methods.text_with_panel.accuracy; paper/data/run11-simulated-readers/simulated-readers-claude-sonnet-5-5-20261106-2dc461f52ea4.json methods.text_with_panel.minus_plain.estimate; paper/data/run11-simulated-readers/simulated-readers-claude-sonnet-5-5-20261106-2dc461f52ea4.json methods.text_with_panel.minus_plain.lo; paper/data/run11-simulated-readers/simulated-readers-claude-sonnet-5-5-20261106-2dc461f52ea4.json methods.text_with_panel.minus_plain.hi; paper/data/run11-simulated-readers/simulated-readers-claude-sonnet-5-5-20261106-2dc461f52ea4.json methods.text_with_panel.minus_always_reference.estimate; paper/data/run11-simulated-readers/simulated-readers-claude-sonnet-5-5-20261106-2dc461f52ea4.json methods.text_with_panel.minus_always_reference.lo; paper/data/run11-simulated-readers/simulated-readers-claude-sonnet-5-5-20261106-2dc461f52ea4.json methods.text_with_panel.minus_always_reference.hi
\quad 5. text classifier with the panel & 0.5484 & 0.0251 (0.0098, 0.0402) & -0.0216 (-0.0367, -0.0060) \\
% src: paper/data/run11-simulated-readers/simulated-readers-claude-sonnet-5-5-20261106-2dc461f52ea4.json methods.examples_in_prompt.accuracy; paper/data/run11-simulated-readers/simulated-readers-claude-sonnet-5-5-20261106-2dc461f52ea4.json methods.examples_in_prompt.minus_plain.estimate; paper/data/run11-simulated-readers/simulated-readers-claude-sonnet-5-5-20261106-2dc461f52ea4.json methods.examples_in_prompt.minus_plain.lo; paper/data/run11-simulated-readers/simulated-readers-claude-sonnet-5-5-20261106-2dc461f52ea4.json methods.examples_in_prompt.minus_plain.hi; paper/data/run11-simulated-readers/simulated-readers-claude-sonnet-5-5-20261106-2dc461f52ea4.json methods.examples_in_prompt.minus_always_reference.estimate; paper/data/run11-simulated-readers/simulated-readers-claude-sonnet-5-5-20261106-2dc461f52ea4.json methods.examples_in_prompt.minus_always_reference.lo; paper/data/run11-simulated-readers/simulated-readers-claude-sonnet-5-5-20261106-2dc461f52ea4.json methods.examples_in_prompt.minus_always_reference.hi
\quad 6. earlier results in the prompt & 0.5419 & 0.0185 (0.0073, 0.0301) & -0.0282 (-0.0487, -0.0102) \\
% src: paper/data/run11-simulated-readers/simulated-readers-claude-sonnet-5-5-20261106-2dc461f52ea4.json methods.audience_brief.accuracy; paper/data/run11-simulated-readers/simulated-readers-claude-sonnet-5-5-20261106-2dc461f52ea4.json methods.audience_brief.minus_plain.estimate; paper/data/run11-simulated-readers/simulated-readers-claude-sonnet-5-5-20261106-2dc461f52ea4.json methods.audience_brief.minus_plain.lo; paper/data/run11-simulated-readers/simulated-readers-claude-sonnet-5-5-20261106-2dc461f52ea4.json methods.audience_brief.minus_plain.hi; paper/data/run11-simulated-readers/simulated-readers-claude-sonnet-5-5-20261106-2dc461f52ea4.json methods.audience_brief.minus_always_reference.estimate; paper/data/run11-simulated-readers/simulated-readers-claude-sonnet-5-5-20261106-2dc461f52ea4.json methods.audience_brief.minus_always_reference.lo; paper/data/run11-simulated-readers/simulated-readers-claude-sonnet-5-5-20261106-2dc461f52ea4.json methods.audience_brief.minus_always_reference.hi
\quad 7. audience brief & 0.5828 & 0.0594 (0.0455, 0.0715) & 0.0127 (-0.0067, 0.0308) \\
\midrule
\multicolumn{4}{@{}l}{\emph{Haiku}} \\
% src: paper/data/run11-simulated-readers/simulated-readers-claude-haiku-5-5-20261106-2dc461f52ea4.json plain.accuracy; paper/data/run11-simulated-readers/simulated-readers-claude-haiku-5-5-20261106-2dc461f52ea4.json readings.plain_minus_always_reference.estimate; paper/data/run11-simulated-readers/simulated-readers-claude-haiku-5-5-20261106-2dc461f52ea4.json readings.plain_minus_always_reference.lo; paper/data/run11-simulated-readers/simulated-readers-claude-haiku-5-5-20261106-2dc461f52ea4.json readings.plain_minus_always_reference.hi
\quad plain panel & 0.5144 &  & -0.0556 (-0.0747, -0.0385) \\
% src: paper/data/run11-simulated-readers/simulated-readers-claude-haiku-5-5-20261106-2dc461f52ea4.json methods.direction_map.accuracy; paper/data/run11-simulated-readers/simulated-readers-claude-haiku-5-5-20261106-2dc461f52ea4.json methods.direction_map.minus_plain.estimate; paper/data/run11-simulated-readers/simulated-readers-claude-haiku-5-5-20261106-2dc461f52ea4.json methods.direction_map.minus_plain.lo; paper/data/run11-simulated-readers/simulated-readers-claude-haiku-5-5-20261106-2dc461f52ea4.json methods.direction_map.minus_plain.hi; paper/data/run11-simulated-readers/simulated-readers-claude-haiku-5-5-20261106-2dc461f52ea4.json methods.direction_map.minus_always_reference.estimate; paper/data/run11-simulated-readers/simulated-readers-claude-haiku-5-5-20261106-2dc461f52ea4.json methods.direction_map.minus_always_reference.lo; paper/data/run11-simulated-readers/simulated-readers-claude-haiku-5-5-20261106-2dc461f52ea4.json methods.direction_map.minus_always_reference.hi
\quad 1. direction map & 0.5669 & 0.0525 (0.0357, 0.0685) & -0.0032 (-0.0117, 0.0042) \\
% src: paper/data/run11-simulated-readers/simulated-readers-claude-haiku-5-5-20261106-2dc461f52ea4.json methods.reader_weights.accuracy; paper/data/run11-simulated-readers/simulated-readers-claude-haiku-5-5-20261106-2dc461f52ea4.json methods.reader_weights.minus_plain.estimate; paper/data/run11-simulated-readers/simulated-readers-claude-haiku-5-5-20261106-2dc461f52ea4.json methods.reader_weights.minus_plain.lo; paper/data/run11-simulated-readers/simulated-readers-claude-haiku-5-5-20261106-2dc461f52ea4.json methods.reader_weights.minus_plain.hi; paper/data/run11-simulated-readers/simulated-readers-claude-haiku-5-5-20261106-2dc461f52ea4.json methods.reader_weights.minus_always_reference.estimate; paper/data/run11-simulated-readers/simulated-readers-claude-haiku-5-5-20261106-2dc461f52ea4.json methods.reader_weights.minus_always_reference.lo; paper/data/run11-simulated-readers/simulated-readers-claude-haiku-5-5-20261106-2dc461f52ea4.json methods.reader_weights.minus_always_reference.hi
\quad 2. reader weights & 0.5375 & 0.0231 (0.0101, 0.0349) & -0.0325 (-0.0531, -0.0145) \\
% src: paper/data/run11-simulated-readers/simulated-readers-claude-haiku-5-5-20261106-2dc461f52ea4.json methods.audience_matching.accuracy; paper/data/run11-simulated-readers/simulated-readers-claude-haiku-5-5-20261106-2dc461f52ea4.json methods.audience_matching.minus_plain.estimate; paper/data/run11-simulated-readers/simulated-readers-claude-haiku-5-5-20261106-2dc461f52ea4.json methods.audience_matching.minus_plain.lo; paper/data/run11-simulated-readers/simulated-readers-claude-haiku-5-5-20261106-2dc461f52ea4.json methods.audience_matching.minus_plain.hi; paper/data/run11-simulated-readers/simulated-readers-claude-haiku-5-5-20261106-2dc461f52ea4.json methods.audience_matching.minus_always_reference.estimate; paper/data/run11-simulated-readers/simulated-readers-claude-haiku-5-5-20261106-2dc461f52ea4.json methods.audience_matching.minus_always_reference.lo; paper/data/run11-simulated-readers/simulated-readers-claude-haiku-5-5-20261106-2dc461f52ea4.json methods.audience_matching.minus_always_reference.hi
\quad 3. audience matching & 0.5020 & -0.0124 (-0.0231, -0.0002) & -0.0681 (-0.0865, -0.0489) \\
% src: paper/data/run11-simulated-readers/simulated-readers-claude-haiku-5-5-20261106-2dc461f52ea4.json methods.error_correction.accuracy; paper/data/run11-simulated-readers/simulated-readers-claude-haiku-5-5-20261106-2dc461f52ea4.json methods.error_correction.minus_plain.estimate; paper/data/run11-simulated-readers/simulated-readers-claude-haiku-5-5-20261106-2dc461f52ea4.json methods.error_correction.minus_plain.lo; paper/data/run11-simulated-readers/simulated-readers-claude-haiku-5-5-20261106-2dc461f52ea4.json methods.error_correction.minus_plain.hi; paper/data/run11-simulated-readers/simulated-readers-claude-haiku-5-5-20261106-2dc461f52ea4.json methods.error_correction.minus_always_reference.estimate; paper/data/run11-simulated-readers/simulated-readers-claude-haiku-5-5-20261106-2dc461f52ea4.json methods.error_correction.minus_always_reference.lo; paper/data/run11-simulated-readers/simulated-readers-claude-haiku-5-5-20261106-2dc461f52ea4.json methods.error_correction.minus_always_reference.hi
\quad 4. error correction & 0.5360 & 0.0216 (0.0129, 0.0298) & -0.0340 (-0.0525, -0.0175) \\
% src: paper/data/run11-simulated-readers/simulated-readers-claude-haiku-5-5-20261106-2dc461f52ea4.json methods.text_with_panel.accuracy; paper/data/run11-simulated-readers/simulated-readers-claude-haiku-5-5-20261106-2dc461f52ea4.json methods.text_with_panel.minus_plain.estimate; paper/data/run11-simulated-readers/simulated-readers-claude-haiku-5-5-20261106-2dc461f52ea4.json methods.text_with_panel.minus_plain.lo; paper/data/run11-simulated-readers/simulated-readers-claude-haiku-5-5-20261106-2dc461f52ea4.json methods.text_with_panel.minus_plain.hi; paper/data/run11-simulated-readers/simulated-readers-claude-haiku-5-5-20261106-2dc461f52ea4.json methods.text_with_panel.minus_always_reference.estimate; paper/data/run11-simulated-readers/simulated-readers-claude-haiku-5-5-20261106-2dc461f52ea4.json methods.text_with_panel.minus_always_reference.lo; paper/data/run11-simulated-readers/simulated-readers-claude-haiku-5-5-20261106-2dc461f52ea4.json methods.text_with_panel.minus_always_reference.hi
\quad 5. text classifier with the panel & 0.5499 & 0.0355 (0.0216, 0.0518) & -0.0201 (-0.0355, -0.0038) \\
% src: paper/data/run11-simulated-readers/simulated-readers-claude-haiku-5-5-20261106-2dc461f52ea4.json methods.examples_in_prompt.accuracy; paper/data/run11-simulated-readers/simulated-readers-claude-haiku-5-5-20261106-2dc461f52ea4.json methods.examples_in_prompt.minus_plain.estimate; paper/data/run11-simulated-readers/simulated-readers-claude-haiku-5-5-20261106-2dc461f52ea4.json methods.examples_in_prompt.minus_plain.lo; paper/data/run11-simulated-readers/simulated-readers-claude-haiku-5-5-20261106-2dc461f52ea4.json methods.examples_in_prompt.minus_plain.hi; paper/data/run11-simulated-readers/simulated-readers-claude-haiku-5-5-20261106-2dc461f52ea4.json methods.examples_in_prompt.minus_always_reference.estimate; paper/data/run11-simulated-readers/simulated-readers-claude-haiku-5-5-20261106-2dc461f52ea4.json methods.examples_in_prompt.minus_always_reference.lo; paper/data/run11-simulated-readers/simulated-readers-claude-haiku-5-5-20261106-2dc461f52ea4.json methods.examples_in_prompt.minus_always_reference.hi
\quad 6. earlier results in the prompt & 0.5189 & 0.0045 (-0.0095, 0.0180) & -0.0511 (-0.0703, -0.0348) \\
% src: paper/data/run11-simulated-readers/simulated-readers-claude-haiku-5-5-20261106-2dc461f52ea4.json methods.audience_brief.accuracy; paper/data/run11-simulated-readers/simulated-readers-claude-haiku-5-5-20261106-2dc461f52ea4.json methods.audience_brief.minus_plain.estimate; paper/data/run11-simulated-readers/simulated-readers-claude-haiku-5-5-20261106-2dc461f52ea4.json methods.audience_brief.minus_plain.lo; paper/data/run11-simulated-readers/simulated-readers-claude-haiku-5-5-20261106-2dc461f52ea4.json methods.audience_brief.minus_plain.hi; paper/data/run11-simulated-readers/simulated-readers-claude-haiku-5-5-20261106-2dc461f52ea4.json methods.audience_brief.minus_always_reference.estimate; paper/data/run11-simulated-readers/simulated-readers-claude-haiku-5-5-20261106-2dc461f52ea4.json methods.audience_brief.minus_always_reference.lo; paper/data/run11-simulated-readers/simulated-readers-claude-haiku-5-5-20261106-2dc461f52ea4.json methods.audience_brief.minus_always_reference.hi
\quad 7. audience brief & 0.5671 & 0.0527 (0.0392, 0.0687) & -0.0029 (-0.0206, 0.0147) \\
\bottomrule
\end{tabular}

\end{table*}

\begin{table}[t]\centering\footnotesize
% src: tables/fmt-readers-transport.tex; paper/data/run12-simulated-readers-api/simulated-readers-transport-claude-opus-5-5-20261106-2dc461f52ea4.json plain.accuracy_api_minus_tool.resamples 500
\caption{The plain panels of the two runs, which send the same requests by two routes. The difference in direction accuracy has a 90 percent bootstrap interval from 500 resamples of blocks and belongs to no family. The share with the same direction is taken over the contrasts on which neither panel is tied.}
\label{tab:fmt-readers-transport}
\resizebox{\columnwidth}{!}{% generated by paper/build_tables.py from the report keys named above each row (readers); it reads paper/data/run12-simulated-readers-api
\begin{tabular}{lrrrrr}
\toprule
model & run 11 & run 12 & run 12 minus run 11 & same direction, & contrasts \\
 & & & & share & compared \\
\midrule
% src: paper/data/run12-simulated-readers-api/simulated-readers-transport-claude-opus-5-5-20261106-2dc461f52ea4.json plain.accuracy_tool; paper/data/run12-simulated-readers-api/simulated-readers-transport-claude-opus-5-5-20261106-2dc461f52ea4.json plain.accuracy_api; paper/data/run12-simulated-readers-api/simulated-readers-transport-claude-opus-5-5-20261106-2dc461f52ea4.json plain.accuracy_api_minus_tool.estimate; paper/data/run12-simulated-readers-api/simulated-readers-transport-claude-opus-5-5-20261106-2dc461f52ea4.json plain.accuracy_api_minus_tool.lo; paper/data/run12-simulated-readers-api/simulated-readers-transport-claude-opus-5-5-20261106-2dc461f52ea4.json plain.accuracy_api_minus_tool.hi; paper/data/run12-simulated-readers-api/simulated-readers-transport-claude-opus-5-5-20261106-2dc461f52ea4.json plain.agree_share; paper/data/run12-simulated-readers-api/simulated-readers-transport-claude-opus-5-5-20261106-2dc461f52ea4.json plain.n
Opus & 0.5742 & 0.5742 & 0.0000 (-0.0085, 0.0081) & 0.852 & 5,962 \\
% src: paper/data/run12-simulated-readers-api/simulated-readers-transport-claude-sonnet-5-5-20261106-2dc461f52ea4.json plain.accuracy_tool; paper/data/run12-simulated-readers-api/simulated-readers-transport-claude-sonnet-5-5-20261106-2dc461f52ea4.json plain.accuracy_api; paper/data/run12-simulated-readers-api/simulated-readers-transport-claude-sonnet-5-5-20261106-2dc461f52ea4.json plain.accuracy_api_minus_tool.estimate; paper/data/run12-simulated-readers-api/simulated-readers-transport-claude-sonnet-5-5-20261106-2dc461f52ea4.json plain.accuracy_api_minus_tool.lo; paper/data/run12-simulated-readers-api/simulated-readers-transport-claude-sonnet-5-5-20261106-2dc461f52ea4.json plain.accuracy_api_minus_tool.hi; paper/data/run12-simulated-readers-api/simulated-readers-transport-claude-sonnet-5-5-20261106-2dc461f52ea4.json plain.agree_share; paper/data/run12-simulated-readers-api/simulated-readers-transport-claude-sonnet-5-5-20261106-2dc461f52ea4.json plain.n
Sonnet & 0.5234 & 0.5279 & 0.0045 (-0.0040, 0.0134) & 0.812 & 5,964 \\
% src: paper/data/run12-simulated-readers-api/simulated-readers-transport-claude-haiku-5-5-20261106-2dc461f52ea4.json plain.accuracy_tool; paper/data/run12-simulated-readers-api/simulated-readers-transport-claude-haiku-5-5-20261106-2dc461f52ea4.json plain.accuracy_api; paper/data/run12-simulated-readers-api/simulated-readers-transport-claude-haiku-5-5-20261106-2dc461f52ea4.json plain.accuracy_api_minus_tool.estimate; paper/data/run12-simulated-readers-api/simulated-readers-transport-claude-haiku-5-5-20261106-2dc461f52ea4.json plain.accuracy_api_minus_tool.lo; paper/data/run12-simulated-readers-api/simulated-readers-transport-claude-haiku-5-5-20261106-2dc461f52ea4.json plain.accuracy_api_minus_tool.hi; paper/data/run12-simulated-readers-api/simulated-readers-transport-claude-haiku-5-5-20261106-2dc461f52ea4.json plain.agree_share; paper/data/run12-simulated-readers-api/simulated-readers-transport-claude-haiku-5-5-20261106-2dc461f52ea4.json plain.n
Haiku & 0.5144 & 0.5223 & 0.0079 (-0.0035, 0.0206) & 0.671 & 5,964 \\
\bottomrule
\end{tabular}
}
\end{table}

% src: tables/fmt-readers-readings.tex (the fixed words of the three readings: not shown to match more often, raised, raised); paper/data/run11-simulated-readers/simulated-readers-claude-opus-5-5-20261106-2dc461f52ea4.json readings.plain_minus_always_reference.verdict contains_zero, readings.weights_minus_plain.verdict above, readings.text_with_panel_minus_text.verdict above; no number
\paragraph{The three readings.} In the words fixed beforehand: on the headline-only tests of the exploratory split, with nothing hidden from the model's training and contamination not excluded, a panel of ten simulated readers is not shown to match the observed direction more often than always choosing the earliest arm. Aligning the readers with earlier results raised that share, and adding the simulated readers to a classifier fitted on earlier tests raised its share of directions matched. Table~\ref{tab:fmt-readers-readings} holds the numbers, and the three paragraphs below say how far each sentence carries.

% src: tables/fmt-readers-readings.tex rows plain panel, Opus 0.5742 and 0.5742, the earliest-arm rule 0.5701, reading 1 Opus 0.0041 (-0.0174, 0.0256) in run 11 and 0.0041 (-0.0171, 0.0253) in run 12, Sonnet -0.0467 (-0.0704, -0.0230) and -0.0422 (-0.0637, -0.0206), Haiku -0.0556 (-0.0791, -0.0322) and -0.0478 (-0.0707, -0.0248)
\emph{Reading 1.} The primary model's plain panel is right on 0.574 of the contrasts, against 0.570 for the earliest-arm rule. The adjusted interval of the difference contains zero in both runs. The plain panels of the two smaller models are below the rule in both runs, with intervals that lie below zero. A panel of simulated readers, used as it comes, is therefore not shown to be better here than a rule that reads nothing, and for two of the three models its accuracy is lower, with adjusted intervals below zero in both runs.

% src: tables/fmt-readers-readings.tex reading 2 Opus 0.0140 (0.0004, 0.0276) in run 11 and 0.0110 (-0.0035, 0.0254) in run 12, Sonnet 0.0365 (0.0222, 0.0507) and 0.0326 (0.0162, 0.0490), Haiku 0.0231 (0.0073, 0.0389) and -0.0012 (-0.0176, 0.0152); tables/fmt-readers-methods.tex rows 2. reader weights, minus the earliest-arm rule: Sonnet -0.0102 (-0.0308, 0.0082), Haiku -0.0325 (-0.0531, -0.0145)
\emph{Reading 2.} For the primary model in run 11, reader weights add 0.0140 to the plain panel's accuracy, and the lower end of the adjusted interval is 0.0004. The fixed word is therefore \emph{raised}. We do not read it as established: in run 12 the same reading gives an adjusted interval that contains zero. For the middle model the gain is larger and its interval lies above zero in both runs, and for the smallest model the gain of run 11 is absent in run 12. In no case do the weights lift a smaller model's panel above the earliest-arm rule (Table~\ref{tab:fmt-readers-methods}).

% src: tables/fmt-readers-readings.tex reading 3 Opus 0.0203 (0.0100, 0.0306) in run 11 and 0.0184 (0.0092, 0.0276) in run 12, Sonnet 0.0000 (-0.0068, 0.0068), Haiku 0.0015 (-0.0045, 0.0075), the text classifier 0.5484, the earliest-arm rule 0.5701; tables/fmt-readers-methods.tex Opus row 5. text classifier with the panel 0.5687, minus the earliest-arm rule -0.0013 (-0.0153, 0.0122)
\emph{Reading 3.} Adding the primary model's panel to the text classifier raises the classifier's accuracy by 0.0203 in run 11 and by 0.0184 in run 12, and both adjusted intervals lie above zero. This is the one reading that both runs support. Its size needs care. The classifier alone is below the earliest-arm rule, and with the panel it reaches 0.569, which is still not above that rule's 0.570. With the panel the classifier is no longer shown to be below the rule that reads nothing, and it is not above it. Nor is the pair shown to differ from the plain panel alone (Table~\ref{tab:fmt-readers-methods}), so the gain says that the classifier is weak here more than that the two sources add up. The panels of the two smaller models add nothing to the classifier that we can detect.

% src: tables/fmt-readers-methods.tex Opus row 7. audience brief 0.6122, minus the earliest-arm rule 0.0422 (0.0260, 0.0575), Sonnet 0.0127 (-0.0067, 0.0308), Haiku -0.0029 (-0.0206, 0.0147); tables/fmt-readers-methods-api.tex Opus row 7. audience brief 0.6170, minus the earliest-arm rule 0.0469 (0.0300, 0.0623), Haiku 0.5401, minus the earliest-arm rule -0.0300 (-0.0472, -0.0131); tables/fmt-readers-readings.tex the direction ceiling 0.8301
\paragraph{The seven methods.} Table~\ref{tab:fmt-readers-methods} sets the methods side by side. One stands out for the primary model: with the audience brief its panel is right on 0.612 of the contrasts in run 11 and on 0.617 in run 12. Its difference from the earliest-arm rule is 0.042 in run 11, with an interval that lies above zero, and run 12 agrees. We report this as a description and not as a finding. The comparison was not one of the three readings, and its interval is not adjusted for the seven methods and three models of the table. The model writes a new note in each run, so for this method run 12 is a second draw and not a repeat of the same requests. The design expected this method to be among the weakest, so what the table shows was not predicted. For the smallest model the same method is below the earliest-arm rule in run 12. Even the best figure here is far from the ceiling of 0.830. Audience matching, the reweighting that uses no outcome of a test, has the lowest or the second lowest accuracy of the seven methods for every model in both runs.

% src: tables/fmt-readers-transport.tex Opus 0.5742, 0.5742, 0.0000 (-0.0085, 0.0081), same direction 0.852; Sonnet 0.0045 (-0.0040, 0.0134); Haiku 0.0079 (-0.0035, 0.0206), same direction 0.671; tables/fmt-readers-more.tex row the two halves of the panel give the same direction, share: Opus 0.750
\paragraph{The two runs against each other.} For the primary model the plain panel's accuracy is the same in the two runs to four decimals, and the interval of the difference runs from minus 0.0085 to 0.0081 (Table~\ref{tab:fmt-readers-transport}). By the rule fixed beforehand, the context that the command line tool adds did not move the plain panel's accuracy by an amount this design can detect. The corresponding interval contains zero for each of the two smaller models. Equal accuracy is not equal answers: the two panels of the primary model give the same direction on 0.852 of the contrasts, and those of the smallest model on 0.671. Within one run, two halves of the primary model's panel agree on 0.750 of the contrasts. A panel of ten readers is thus a noisy instrument, and a difference between two runs that lies inside the intervals of Table~\ref{tab:fmt-readers-transport} should not be read.

% src: tables/fmt-readers-probes.tex (completion: winner minus other, each within two standard errors of zero in both runs; familiarity: 0.509 to 0.534; named publisher: 0.557 to 0.595)
\paragraph{Earlier exposure to the archive.} Three probes ask whether a model recalls the winning headline of a test better than a losing one, which it could if the publisher ran the winner afterwards (Table~\ref{tab:fmt-readers-probes} in Appendix~\ref{sec:readers-detail}). Given the first four words of a headline, a model completes the winning headline of a test no better than another headline of the same test. Its ratings of how familiar a headline is point to the winner on 0.509 to 0.534 of the contrasts probed. The reports give no interval for this share, and it is above one half for every model in both runs. Asked directly which headline of a named publisher's pair won, each model is right more often than a coin would be, which it can do by judging the headlines and which is therefore no evidence either way. Two probes thus give no clear sign that winners are recalled better, and the third cannot tell. None of them can rule recall out, and every statement of this section is made with contamination not excluded.

% src: sections/limitations.tex, the paragraph on the simulated readers; no number
\paragraph{What this section does not show.} The result is one measurement on one archive with assumed profiles, one vendor and no fresh split. A null here is as hard to read as a gain: the profiles, the feed format or the models could each explain it. Agreement between the three models is not confirmation, because they share a vendor and, most likely, training data. Section~\ref{sec:limitations} lists these limits with the others.
